\documentclass[10pt]{article}

% ---------- IDIOMA Y CODIFICACIÓN ----------
\usepackage[utf8]{inputenc}
\usepackage[spanish]{babel}
\usepackage[margin=2.5cm]{geometry}
\usepackage{setspace}
\usepackage{graphicx}

% ---------- TABLAS (Actividad 1: identificación de dispositivos) ----------
\usepackage{booktabs}   % líneas de tabla profesionales (\toprule, \midrule, \bottomrule)
\usepackage{array}
\usepackage{ragged2e}
\usepackage[table]{xcolor} % opcional, para color en encabezado
\usepackage{tabularx}   % columnas que se ajustan al ancho de página
\usepackage{longtable}  % tablas que pueden partirse entre páginas (útil si tu tabla de dispositivos es larga)
\usepackage{caption}    % control de estilo de "Tabla X: ..."

% ---------- GLOSARIO (Actividad 2) ----------
% Enfoque simple con "description", sin necesidad de correr makeglossaries aparte
% (evita agregar un paso extra a tu pipeline de compilación).
\usepackage{enumitem}
\newlist{glosario}{description}{1}
\setlist[glosario]{
    font=\normalfont\bfseries,
    leftmargin=0pt,
    labelsep=0.5em,
    itemsep=0.5em
}
% Uso dentro del documento:
%   \begin{glosario}
%       \item[Chipset] Definición y características...
%   \end{glosario}

% ---------- PLACEHOLDER \relleno ----------
% Uso: \relleno              -> texto de aviso genérico
%      \relleno{Tu texto}    -> muestra tu texto en rojo (útil para ir dejando borradores)
\newcommand{\relleno}[1][\textbf{[COMPLETAR AQUÍ]}]{\textcolor{red}{#1}}
% ---------- CONFIGURACIONES PARA ENCABEZADO Y PORTADA ----------
\usepackage{colortbl}
\usepackage{float}

% Definición de colores y columnas personalizadas (usadas en la portada)
\definecolor{tablebackground}{rgb}{0.8, 0, 0}
\newcolumntype{x}[1]{>{\centering\arraybackslash\hspace{0pt}}p{#1}}
\newcolumntype{M}[1]{>{\centering\arraybackslash}m{#1}}
\newcolumntype{N}{@{}m{0pt}@{}}
\newcolumntype{P}[1]{>{\RaggedRight\arraybackslash}p{#1}}

% Variables de datos para la portada y encabezado
\newcommand{\itemEmail}{ccapatinta@unsa.edu.pe}
\newcommand{\itemStudent}{Camilo B.A. Capatinta Almanza (Teoría grupo B)}
\newcommand{\itemCourse}{Lab. Arq. de computadoras}
\newcommand{\itemCourseCode}{1702227}
\newcommand{\itemSemester}{IV}
\newcommand{\itemUniversity}{Universidad Nacional de San Agustín de Arequipa}
\newcommand{\itemFaculty}{Facultad de Ingeniería de Producción y Servicios}
\newcommand{\itemDepartment}{Departamento Académico de Ingeniería de Sistemas e Informática}
\newcommand{\itemSchool}{Escuela Profesional de Ingeniería de Sistemas}
\newcommand{\itemAcademic}{2026 - B}
\newcommand{\itemInput}{15/09/2026}
\newcommand{\itemOutput}{22/09/2026}
\newcommand{\itemPracticeNumber}{02}
\newcommand{\itemTheme}{Sistemas Digitales}

% Configuración del encabezado y pie de página (fancyhdr)
\usepackage{fancyhdr}
\pagestyle{fancy}
\fancyhf{}
\setlength{\headheight}{40.51407pt}
\renewcommand{\headrulewidth}{1pt}
\renewcommand{\footrulewidth}{1pt}

% Encabezado
\fancyhead[L]{\raisebox{-0.2\height}{\includegraphics[width=3cm]{img/logo_episunsa.png}}}
\fancyhead[C]{\fontsize{7}{7}\selectfont    \itemUniversity \\ \itemFaculty \\ \itemDepartment \\ \itemSchool \\ \textbf{\itemCourse}}
\fancyhead[R]{\raisebox{-0.2\height}{\includegraphics[width=1.2cm]{img/logo_abet.png}}}

% Pie de página
\fancyfoot[L]{Estudiante Camilo Capatinta}
\fancyfoot[C]{\itemCourse}
\fancyfoot[R]{Página \thepage}

\begin{document}

% --- INICIO DE LA PORTADA ---
\vspace*{10px}
    
\begin{center}  
    \fontsize{17}{17} \textbf{ Informe de Laboratorio \itemPracticeNumber}
\end{center}
\centerline{\textbf{\Large Tema: \itemTheme}}

\begin{flushright}
    \begin{tabular}{|M{2.5cm}|N|}
        \hline 
        \rowcolor{tablebackground}
        \color{white} \textbf{Nota}  \\
        \hline 
             \\[30pt]
        \hline          
    \end{tabular}
\end{flushright}    

\begin{table}[h!]
    \begin{tabular}{|x{4.7cm}|x{4.8cm}|x{4.8cm}|}        
        \hline 
        \rowcolor{tablebackground}
        \color{white} \textbf{Estudiante} & \color{white}\textbf{Escuela}  & \color{white}\textbf{Asignatura}   \\
        \hline 
        {\itemStudent \par \itemEmail} & \itemSchool & {\itemCourse \par Semestre: \itemSemester \par Código: \itemCourseCode}     \\
        \hline          
    \end{tabular}
\end{table}     

\begin{table}[H]
    \begin{tabular}{|x{4.7cm}|x{4.8cm}|x{4.8cm}|}
        \hline 
        \rowcolor{tablebackground}
        \color{white}\textbf{Laboratorio} & \color{white}\textbf{Tema}  & \color{white}\textbf{Duración}   \\
        \hline 
        \itemPracticeNumber & \itemTheme & 02 horas   \\
        \hline 
    \end{tabular}
\end{table}

\begin{table}[H]
    \begin{tabular}{|x{4.7cm}|x{4.8cm}|x{4.8cm}|}
        \hline 
        \rowcolor{tablebackground}
        \color{white}\textbf{Semestre académico} & \color{white}\textbf{Fecha de inicio}  & \color{white}\textbf{Fecha de entrega}   \\
        \hline 
        \itemAcademic & \itemInput &  \itemOutput  \\
        \hline 
    \end{tabular}
\end{table}

\newpage
% --- FIN DE LA PORTADA ---

\section{Actividades}

\begin{itemize}
    \item Diseñar e implementar un sumador total de un bit, con acarreo de entrada y salida.
    \item Desarrollar el diseño y la implementación de un sumador total de un bit (con acarreo de entrada y salida). El diseño debe contener la tabla de verdad, los mini-términos, las funciones lógicas para la suma y el acarreo de salida.
    \item Luego de obtener las funciones lógicas debe desarrollar las simplificaciones y finalmente elaborar el circuito lógico.
    \item El circuito lógico se debe implementar en el simulador de circuitos digitales o similares.
    \item Finalmente, demostrar el diseño y el funcionamiento del sistema digital denominado sumador total.
\end{itemize}

\section{Conclusiones}

\end{document}